\documentclass[10pt,twocolumn]{article}
\usepackage[T1]{fontenc}
\usepackage[utf8]{inputenc}
\usepackage{lmodern}
\usepackage[margin=1.8cm,top=2.4cm,bottom=2cm,columnsep=0.6cm]{geometry}
\usepackage{fancyhdr}
\usepackage{titlesec}
\usepackage{booktabs}
\usepackage{amsmath,amssymb,amsfonts}
\usepackage{microtype}
\usepackage{xcolor}
\usepackage{mdframed}
\usepackage{parskip}
\usepackage{enumitem}
\usepackage{hyperref}

\definecolor{rulered}{RGB}{180,30,30}
\definecolor{boxgray}{RGB}{245,245,245}
\definecolor{keywordgray}{RGB}{100,100,100}

\pagestyle{fancy}
\fancyhf{}
\fancyhead[L]{\textsc{\small Claude Mag}}
\fancyhead[C]{\small\textit{Machine Learning Research Digest}}
\fancyhead[R]{\small 2026-08-05}
\fancyfoot[C]{\small\thepage}
\renewcommand{\headrulewidth}{0.8pt}

\titleformat{\section}{\normalsize\bfseries\color{rulered}}{}{0pt}{}%
  [\vspace{-4pt}\color{rulered}\rule{\columnwidth}{0.4pt}]
\titlespacing{\section}{0pt}{8pt}{4pt}

\hypersetup{colorlinks=true,urlcolor=rulered,linkcolor=black}

\begin{document}
\setlength{\parindent}{0pt}
\setlength{\parskip}{4pt}

{\small\color{keywordgray}\textbf{Keywords:}
  tabular prediction \textbullet{}
  in-context learning \textbullet{}
  dimensionality \textbullet{}
  LLM limitations}\par\vspace{4pt}

{\LARGE\bfseries\raggedright
  Why Large Language Models Fail at Tabular Prediction}\par\vspace{4pt}

{\large\itshape\raggedright
  A Systematic Study Identifies Input Dimensionality---Not Format, Noise, or
  Tokenisation---as the Decisive Factor in LLM Tabular Failure}\par\vspace{6pt}

{\small\textbf{By} Marta Garnelo \& Wojciech M. Czarnecki
  (Google DeepMind)}\par\vspace{2pt}

{\small\color{keywordgray}arXiv:2608.02412 \textbullet{} Submitted 3 August 2026}
\par\vspace{2pt}\noindent{\color{rulered}\rule{\columnwidth}{0.6pt}}\vspace{6pt}

Large language models are outperformed by fifty-year-old baselines on tabular
data. A new controlled study across 31 benchmark datasets falsifies four popular
explanations for this failure and finds that dimensionality alone is decisive:
LLM accuracy monotonically falls as feature count grows, while every classical
method stays flat or improves, and in two dimensions the LLM's predictions match
a local distance-based method with up to 91.6\% grid agreement.

\begin{mdframed}[backgroundcolor=boxgray,linecolor=rulered,linewidth=1pt,
    innerleftmargin=6pt,innerrightmargin=6pt,
    innertopmargin=6pt,innerbottommargin=6pt]
\textbf{\small AT A GLANCE}\par\vspace{4pt}
\begin{description}[leftmargin=0pt,labelsep=4pt,itemsep=2pt,parsep=0pt,
    font=\small\bfseries]
  \item[Problem:] LLMs routinely lose to classical baselines on tabular
    prediction, but the root cause has been unknown.
  \item[Why it matters:] Tabular data dominates enterprise ML; understanding why
    LLMs fail guides where not to deploy them and points to what a fix would require.
  \item[Method:] Evaluate a frontier LLM in pure in-context inference on 31 datasets;
    sweep input dimensionality via random linear projections; test five candidate hypotheses.
  \item[Result:] Dimensionality is the sole decisive factor; all four other hypotheses
    are falsified by controlled experiments.
  \item[Evidence:] 31 datasets, 9 comparison methods; 91.6\% prediction grid agreement
    with local distance-based methods in 2-D; failure mode unmatched by any classical degradation.
\end{description}
\end{mdframed}

\section{The Big Picture}

Tabular data---rows of numeric and categorical features with a prediction target---is
the most common format in real-world machine learning. Despite LLMs' remarkable
across-the-board progress, they have failed to close the gap with classical methods
on tabular benchmarks. A frontier LLM prompted with a CSV of training examples and
asked to predict the label for a test row is consistently beaten by random forests,
gradient boosting, and even simple $k$-nearest-neighbour classifiers.

This failure has spawned a dedicated field of tabular foundation models, yet the
underlying cause has been debated without controlled evidence. Garnelo and Czarnecki
run the most systematic study to date, falsifying four popular hypotheses and
identifying the sole culprit: input dimensionality.

\section{Why Existing Approaches Fall Short}

The field has proposed several ad-hoc explanations:
\begin{itemize}[itemsep=1pt,parsep=0pt]
  \item \textbf{Noisy / non-separable data:} LLMs may be sensitive to label noise
    or overlapping class boundaries more common in tabular data than text.
  \item \textbf{CSV format:} Serialising rows as comma-separated values may obscure
    the column structure that classical methods exploit.
  \item \textbf{Numeric tokenisation:} Numbers like \texttt{3.14} are split into
    multiple tokens (\texttt{3}, \texttt{.}, \texttt{14}), possibly disrupting
    numeric reasoning.
  \item \textbf{Query batch size:} Classifying many test points per prompt may
    introduce interference between predictions.
\end{itemize}

None of these has been rigorously tested as a \emph{cause} of the failure.
All four are falsified below.

\section{Core Mechanism}

The study places the LLM in its purest in-context regime: a single generation
pass over a prompt containing the full training set serialised as CSV plus one
test point. No tools, fine-tuning, chain-of-thought, or agentic scaffolding.
This isolates the LLM's native prediction mechanism.

To control dimensionality, random linear projections are applied to each of
the 31 datasets, producing versions with 2, 4, 8, 16, 32, and 64 features.
This lets the study vary dimensionality while keeping the underlying dataset
structure the same.

The LLM is compared to eight classical methods including $k$-NN, random forests,
gradient-boosted trees, and linear classifiers---nine methods total.

\section{Technical Formulation}

Let $\mathbf{X} \in \mathbb{R}^{n \times d}$ be the training set with labels
$\mathbf{y} \in \{1,\ldots,C\}^n$, and let $\mathbf{x}^* \in \mathbb{R}^d$ be
a test point. The LLM predicts:
\[
\hat{y}^* = \arg\max_{c}\, p_{\mathrm{LLM}}(c \mid \mathrm{prompt}(\mathbf{X}, \mathbf{y}, \mathbf{x}^*)).
\]
The dimensionality sweep uses a random matrix $\mathbf{P} \in \mathbb{R}^{d \times k}$
with $k < d$ to project features: $\tilde{\mathbf{x}} = \mathbf{x}\mathbf{P}$.
Accuracy is measured as a function of $k$.

For the behavioural comparison, the \emph{prediction grid agreement} between
two classifiers $f_1$ and $f_2$ is:
\[
\mathrm{Agreement}(f_1, f_2)
= \mathbb{E}_{\mathbf{x}^*}\!\left[\mathbf{1}\!\left[f_1(\mathbf{x}^*) = f_2(\mathbf{x}^*)\right]\right],
\]
evaluated on a fine grid over the feature space.

\section{Learning or Inference Procedure}

The experimental protocol has three phases:

\begin{enumerate}[itemsep=2pt,parsep=0pt]
  \item \textbf{Hypothesis testing:} For each of the five hypotheses, design a
    controlled variant of the prompt or data (e.g., corrupt labels for (a), use
    alternative serialisation for (b)) and test whether it changes the gap to classical
    methods. No change = hypothesis falsified.
  \item \textbf{Dimensionality sweep:} Apply random projections to sweep $d \in
    \{2, 4, 8, 16, 32, 64\}$ and measure accuracy vs.\ $d$ for all nine methods.
  \item \textbf{Behavioural comparison:} In 2-D (where the LLM still performs
    competitively), compute prediction grid agreement between the LLM and all
    eight classical methods to identify which method's decision boundary it mimics.
\end{enumerate}

\section{What the Guarantee Says}

Falsification results are definitive at the experimental level: each variant
produces no statistically significant change in the LLM--classical gap when
hypotheses (a)--(d) are tested. The dimensionality effect is monotone across all
31 datasets and all 9 comparison methods, making it robust to dataset-specific
confounds.

\section{Experimental Findings}

\textbf{Dimensionality is decisive.}
The LLM is the only method among nine whose accuracy decreases monotonically with
$d$. Every classical baseline stays flat or improves as $d$ grows (more features
contain more signal for classical methods). The crossover point where classical
methods overtake the LLM is typically around $d = 4$--8.

\textbf{In 2-D, the LLM mimics a distance-based method.}
Grid agreement between the LLM and the nearest-neighbour classifier reaches
\textbf{91.6\%} in 2-D. No other classical method achieves this level of
agreement with the LLM's predictions at any dimensionality.

\textbf{High-dimensional failure mode is novel.}
At $d > 8$, the LLM's decision boundary no longer matches any classical method,
including noise-corrupted or degraded variants. The failure is mechanistically
distinct from all classical failure modes.

\section{Ablations and Interpretation}

\textbf{Hypothesis (a)---noise:} Adding label noise to all methods equally
reduces accuracy for all, but the LLM--classical gap is unchanged. Falsified.

\textbf{Hypothesis (b)---CSV format:} Alternative serialisations (JSON, Markdown
tables, space-separated) produce no statistically significant change in LLM accuracy.
Falsified.

\textbf{Hypothesis (c)---tokenisation:} Replacing all numeric values with symbolic
tokens (\texttt{A}, \texttt{B}, \ldots) does not restore LLM performance.
Falsified.

\textbf{Hypothesis (d)---batch size:} Classifying 1, 5, or 20 test points per
query produces no consistent change in accuracy. Falsified.

\textbf{The prediction mechanism at high $d$ is an open question.}
The LLM's unique sensitivity to dimensionality and the breakdown of its resemblance
to any known classical method at $d > 8$ leave its high-dimensional prediction
mechanism as a fundamental open problem for mechanistic interpretability research.

\section*{Reference}

Marta Garnelo, Wojciech M. Czarnecki. \textbf{Why Large Language Models Fail at
Tabular Prediction}. arXiv:2608.02412, August 2026.
\url{https://arxiv.org/abs/2608.02412}

\end{document}
